\section{Conclusion}

Knowledge arbitration should be treated as an inference problem, not as a
retrieval heuristic. Once written as a posterior-predictive decision problem,
the right comparison is not between one tuned adaptive rule and another, but
between reliability-aware arbitration and fixed trust policies that cannot
adapt to conflict families. Our corrected benchmark waves show that the
Bayes-style rule wins that comparison cleanly, while fixed trust in retrieval
or fixed trust in parametric memory incurs enormous regret.

The reasoning-time result is also now clearer than it was at the start of the
project. The original monotone claim does not survive contact with real traces.
What survives is a more interesting law: reasoning amplifies overconfidence
under conflict, and the shape of that amplification depends on the conflict
family. WikiContradict exhibits an intermediate-CoT peak with partial recovery.
ConflictBank conflict becomes even harsher at 14B, remaining severely
overconfident through long CoT. The completed same-family Qwen sweep shows that
this should not be oversold as a universal 7B-vs-14B asymmetry: Qwen does not
recover on ConflictBank through 32B, while 32B is the first scale at which
recovery reappears on WikiContradict. The finished theorem-3 claim is therefore
benchmark-dependent two-regime behavior, not a family-invariant scale law. The
cleanest interpretation is an RLVR-conditioned misspecification story:
reasoning-time confidence can improve in low-conflict regimes and collapse in
controlled conflict regimes where endogenous evidence keeps sharpening the wrong
answer.

The next steps are clearer now than the original version suggested. The theory
side still needs fuller writeups for the Bayes-optimal rule and the fixed-policy
minimax lower bound, and the theorem-3 side should continue to be written as a
structural misspecification account rather than a theorem of equal formal
status. Empirically, the remaining high-value gaps are direct RLCR-style
training comparisons, broader checkpoint-family sweeps, and mitigation
experiments such as bagged reasoning and confidence decoupling. But the main
paper arc is no longer speculative: fixed trust policies are wrong, a
reliability-aware rule does better, and real reasoning traces reveal a
structured calibration pathology under knowledge conflict.

\paragraph{Scope and limitations.}
Three practical boundaries matter for interpretation. First, the current paper
does not yet report a completed head-to-head number against RLCR
\citep{damani2025rlcr}; our training-time evidence is still lighter-weight than
a full reward-level redesign on the same backbone. Second, the retrieval demo
is deployment-motivated but not production-scale: it uses bounded
Wikipedia-backed retrieval rather than a continuously refreshed live stack.
Third, the benchmark, prompt, and calibration suite here is English-only.
